\chapter{Conclusions and Future Work} \label{Chap5}

This study addressed a central learner-representation question: how should learner characteristics be represented in LLM-based student simulation when they are expected to influence information processing, knowledge formation, and downstream behaviour through the learning process? Using ADHD-related Attention and Working-Memory characteristics as a controlled case, the thesis compared persona prompting with a Cognitive-Process-Based (CPB) representation and reached a qualified affirmative answer to the Main Research Question. The experiments showed that CPB could translate predefined processing constraints into attributable information- and knowledge-state changes and produce task-sensitive behavioural patterns aligned with the expected directions of controlled distraction and instructional processing demand. In contrast, persona-level differences did not consistently translate into process-sensitive effects of the same kind. After Language Ability and Big-Five response-stage prompts were introduced, specific CPB assessment outcomes changed, but the existing graded constraint structure and associated aggregate patterns remained observable.

The principal contribution of this thesis is a \textbf{functionally grounded learner representation}. It separates selected process-relevant characteristics from persona instructions that directly prescribe how a student should appear, allowing those characteristics instead to alter learner state through constraints on the learning process and allowing the resulting state to influence downstream behaviour. Building on this representation, the thesis also developed a progressive validation framework aligned with the scope of its claims, moving from mechanism execution and information-state transitions to task-sensitive behaviour and attribute-extension analysis. The explicit knowledge state and process traceability therefore serve as methodological advantages of the process-based representation: they make the relationship between behavioural outcomes and the processes that produced them directly inspectable.

These conclusions remain subject to clear scope limitations. Attention and Working Memory are limited functional abstractions of ADHD cognition and do not represent a complete or clinically realistic ADHD learning process. PDB is a surprisal-based proxy for instructional processing demand rather than a direct measure of human cognitive load. The study is also limited to English finance materials, multi-round teaching, and short-answer assessment, and no real ADHD learner data were available to establish external behavioural validity. Future work should therefore validate distraction- and processing-demand-related patterns against data from real ADHD and NT learners, refine probabilistic attention allocation, partial retention, interference, and dynamic processing mechanisms, and extend the evaluation across subject domains, task formats, and base models.

Overall, this thesis does not claim that CPB should replace persona prompting in every setting. Its broader implication is that \textbf{learner representation should match the functional level at which a learner characteristic is expected to operate}. Persona prompting may be sufficient for characteristics that primarily affect linguistic expression or interaction style. For characteristics expected to alter knowledge acquisition and subsequent task behaviour through the learning process, representing them at the corresponding processing layer and allowing behavioural differences to emerge from the resulting learner state provides a direction for developing more modular and functionally grounded LLM-based simulated learners.
